% ============================================================
% CONCLUSION — slides 27-29
% ============================================================

\section{Conclusion}

% ----------------------------------------------------------
% SLIDE 27: Summary of results
% SAY: "Three parts, three levels of generality. Part I gives
%       the geometric law: all spatial variation from body shape.
%       Part II adds exact polarisation and frequency extension.
%       Part III handles coherent MIMO with a closed-form
%       exposure operator. The speedup is six orders of magnitude
%       throughout. Five limitations are irreducible."
% ----------------------------------------------------------
\begin{frame}{Summary of results}

  \begin{columns}[T]
    \begin{column}{0.3\textwidth}
      \begin{center}
        {\small\bfseries\color{navyblue} Part I: mmWave}
      \end{center}
      {\footnotesize
      \[
        \Sab \approx \Sinc\, T_0\, \ReLU(\mu)
      \]
      $10^6\times$ speedup\\
      5--15\,\% error\\
      $\langle P_{\mathrm{abs}}\rangle = \Sinc T_0 \Aab/4$\\[2pt]
      Compliance: $\Sinc < 0.16\,m/(T_0 \Aab)$\\[2pt]
      ReLU network formulation
      }
    \end{column}

    \begin{column}{0.3\textwidth}
      \begin{center}
        {\small\bfseries\color{navyblue} Part II: Generalisations}
      \end{center}
      {\footnotesize
      \[
        P_{\mathrm{abs}} = \mm(\khat) \cdot \ssv_{\mathrm{inc}}
      \]
      Polarisation: $\le 16\,\%$ on body\\
      Multipath ($N\ge20$): $\le 2.5\,\%$\\[2pt]
      $\langle P_{\mathrm{abs}}\rangle = \Sinc \Tbar \Aab/4$\\[2pt]
      Exact \SI{100}{\mega\hertz}--\SI{100}{\giga\hertz}
      }
    \end{column}

    \begin{column}{0.3\textwidth}
      \begin{center}
        {\small\bfseries\color{navyblue} Part III: Coherent MIMO}
      \end{center}
      {\footnotesize
      \[
        \Sab = \|\Gtilde\,\xx\|^2
      \]
      Exposure operator $\QQ = \int \Gtilde^H \Gtilde\,\diff A$\\[2pt]
      Alignment $\rho \in [0,1]$\\[2pt]
      QCQP: $\xx^\star \propto (\lambda\QQ+\nu\II)^{-1}\hh^*$\\[2pt]
      Approx.\ error $\lesssim 5\,\%$
      }
    \end{column}
  \end{columns}

  \vspace{8pt}
  Five irreducible limitations:
  {\footnotesize
  reactive near-field ($d < \lambda/2\pi$) \,$\mid$\,
  coherent multi-bounce \,$\mid$\,
  diffraction at edges \,$\mid$\,
  volumetric SAR below \SI{6}{\giga\hertz} \,$\mid$\,
  tissue dielectric uncertainy (${\sim}20\,\%$)
  }
\end{frame}

% ----------------------------------------------------------
% SLIDE 28: Outlook and dissemination
% SAY: "ArXiv is the most important single action. It establishes
%       priority on everything. Paper A targets TAP or T-EMC:
%       the base framework. It is mostly a condensation exercise.
%       Paper B targets TWC or JSAC: the coherent extension.
%       It needs the simulation study with Sionna RT."
% ----------------------------------------------------------
\begin{frame}{Outlook and dissemination}

  \begin{columns}[T]
    \begin{column}{0.55\textwidth}
      Publication plan:

      \vspace{6pt}
      \begin{enumerate}
        \item \textbf{arXiv} (priority claim)\\
              {\footnotesize Post full monograph. Establishes priority on all results.}

        \vspace{4pt}
        \item \textbf{Paper A} $\rightarrow$ IEEE TAP or T-EMC\\
              {\footnotesize Base framework (Parts I--II).\\
              Mostly condensation from the monograph.}

        \vspace{4pt}
        \item \textbf{Paper B} $\rightarrow$ IEEE TWC or JSAC\\
              {\footnotesize Coherent extension (Part III).\\
              Needs simulation study: Sionna RT + Thelonious phantom.}

        \vspace{4pt}
        \item \textbf{BioEM 2027} abstract\\
              {\footnotesize Reuse existing draft. Part I focus.}
      \end{enumerate}
    \end{column}

    \begin{column}{0.4\textwidth}
      Framing strategy:

      \vspace{6pt}
      {\footnotesize
      Paper A opening:\\
      ``Kodera \textit{et al.}\ (2024) observed numerically that
       ACS converges to projected area above \SI{10}{\giga\hertz}.
       We derive this analytically, identify pseudo-Brewster as
       the physical origin, and extend to local fields, arbitrary
       polarisation, diffuse fields, and sub-GHz frequencies.''
      }

      \vspace{8pt}
      Missing piece for Paper B:\\
      {\footnotesize
        Sionna RT simulation study:\\
        64-element BS, \SI{28}{\giga\hertz},\\
        MRT vs optimal QCQP,\\
        hotspot intensity and SNR penalty.
      }
    \end{column}
  \end{columns}
\end{frame}

% ----------------------------------------------------------
% SLIDE 29: Thank you
% ----------------------------------------------------------
\begin{frame}

  \vfill
  \begin{center}
    {\Large\color{navyblue}\bfseries Thank you}\\[12pt]
    {\normalsize Questions and discussion}\\[16pt]
    {\footnotesize Robin Wydaeghe \quad$\mid$\quad robin.wydaeghe@ugent.be}
  \end{center}
  \vfill
\end{frame}
